\subsection{The tanh network model and quantitative bounds}\label{app:tanhmodel}
Let \(n\ge2\), \(D\ge1\), and \(b\ge1\). A network has activations
\begin{equation}\label{eq:network}
h_{0,j}=x_j\in\{-1,1\},\qquad
h_{\ell,i}=\tanh\!\left(b_{\ell,i}
+\sum_{j=1}^n w_{\ell,ij}h_{\ell-1,j}\right),
\quad 1\le\ell\le D.
\end{equation}
All parameters are multiples of \(2^{-b}\), and
\[
|b_{\ell,i}|\le1,\qquad
S:=\max_{\ell,i}\sum_j|w_{\ell,ij}|\le s.
\]
The exact row norm \(S\) is computed from the weights. The target
bit has probability \((1+h_{D,1}(x))/2\).

Preprocessing, probes, and online work are charged as in
Section~\ref{sec:models}, including every precision refinement.
A \emph{recursive call} is an invocation of the
normalized sampler for a neuron or input, including a repeated
invocation of the same coordinate. The root is counted as well.
Write
\begin{equation}\label{eq:bitlength}
\bits=b+\lceil\log_2(n+D+2)\rceil+16.
\end{equation}
This is at most the bit scale \(B\) of \eqref{eq:main-B}, so every
bound stated here with \(\bits\) also holds with \(B\). A plain \(B\)
in a lemma of this appendix bounds the encoding lengths of the
parameters and addresses that the lemma uses.

\begin{theorem}[Exact sampling upper bound]\label{thm:upper}
One uniform algorithm samples the target exactly for every network
with \(S\le2\), and terminates almost surely. Deterministic preprocessing
uses \(O(Dn^2\bits^6)\) bit operations and \(O(Dn^2\bits)\) stored bits.
Uniformly over all contexts, with \(\delta=(S-1)_+\),
\begin{equation}\label{eq:maincall}
\E N_{\rm calls}\le80(D+1)^2e^{51\delta D}.
\end{equation}
Expected online bit work is
\begin{equation}\label{eq:mainbit}
O\!\left(\bits^4(D+1)^2e^{51\delta D}\right).
\end{equation}
The constants in these \(O\)-bounds are absolute.
\end{theorem}

The explicit power \(2\) in \eqref{eq:maincall} counts recursive
calls. Equation \eqref{eq:mainbit} also charges the manipulation of
finite parameters and refinable scalar probabilities. A stable
fixed-precision coefficient recurrence proves the fourth-power
overhead in Section~\ref{sec:fixedprecision}.
For \(S\le1\), Theorem~\ref{thm:fusedcritical} lowers the power of
depth in \eqref{eq:maincall} and \eqref{eq:mainbit} to \(35/32\),
and Theorem~\ref{thm:zero-bias-linear} gives a linear bound when all
biases are zero.

Our lower bounds hold even with unlimited preprocessing. Set
\[
m=2^{\lfloor\log_2n\rfloor},\qquad
F_{a,0}(z)=z,\qquad F_{a,d+1}(z)=\tanh(aF_{a,d}(z)).
\]
If \(a/m\) is on the parameter grid, a legal network realizes
\(F_{a,D}(m^{-1}\sum_{j=1}^m x_j)\): the first distinguished row
has weights \(a/m\), and each later distinguished row has one
incoming weight \(a\). All other neurons and all biases are zero.

\begin{theorem}[Universal query lower bounds]\label{thm:lower}
For this network and \(a>0\), every exact sampler satisfies
\begin{equation}\label{eq:lowerexact}
\sup_x\E[Q\mid x]\ge
\frac{a^{2D}}{1+(3m-2)a^{2D}H_D/m^2},
\qquad H_D=\sum_{j=0}^{D-1}a^{-2j},
\end{equation}
where \(Q\) counts distinct input probes. If \(a>1\), then
\begin{equation}\label{eq:lowersuper}
\sup_x\E[Q\mid x]\ge
\frac{a^2-1}{4a^2-1}\min\{a^{2D},m\}.
\end{equation}
At \(a=1\) and \(1\le D\le m\),
\begin{equation}\label{eq:lowercritical}
\sup_x\E[Q\mid x]\ge\frac{\sqrt D}{100}.
\end{equation}
These bounds apply to every almost surely terminating exact algorithm.
\end{theorem}

\begin{remark}[Other averaging widths]\label{rem:lower-general-m}
The proof in Section~\ref{sec:lower} uses \(m\) only through the
output \(F_{a,D}(m^{-1}\sum_{j=1}^m x_j)\) and the condition
\(D\le m\) in \eqref{eq:lowercritical}. Theorem~\ref{thm:lower}
therefore holds for every integer \(1\le m\le n\), for example every
power of two \(m\le n\), whenever \(a/m\) is on the parameter grid.
It also holds for every legal network with this output, whatever its
other weights.
\end{remark}

Specialize now to \(L=\lceil\log_2n\rceil\), \(D=L\), and \(b=20L\).
Let \(\cost_n(s)\) be the infimum of worst-network, worst-context
expected online bit work among uniform exact samplers with
\(O(Ln^2\log^6n)\) preprocessing time and storage. The norm \(s\)
is fixed as \(n\) grows.

\begin{corollary}[The norm threshold]\label{cor:threshold}
For every fixed \(0\le s\le1\), \(\cost_n(s)=\log^{O(1)}n\).
For every fixed \(s>1\),
\begin{equation}\label{eq:dichotomy}
\cost_n(s)=n^{\Theta(\min\{\log s,1\})}.
\end{equation}
Explicitly, for \(1<s\le2\),
\begin{equation}\label{eq:explicitexponents}
\Omega_s\!\left(n^{\min\{2\log_2s,1\}}\right)
\le\cost_n(s)\le
\widetilde O_s\!\left(\min\{n^2,n^{51(s-1)/\log2}\}\right).
\end{equation}
For larger fixed \(s\), \(\widetilde O_s(n^2)\) is an upper bound.
\end{corollary}

\begin{corollary}[The transition window]\label{cor:window}
Let \(s_n=1+\delta_n\), \(0\le\delta_n\le1\), with \(D=L\) and
\(b=20L\). If \(\delta_nL=O(\log L)\), the upper bound is
polylogarithmic. If \(\delta_nL/\log L\to\infty\), no polylogarithmic
expected online bound is possible. Quantitatively, if
\(\delta L\ge4\log(30L)\), then for all sufficiently large \(L\)
some legal network and context require at least \(e^{\delta L/4}\)
expected input probes.
\end{corollary}

\begin{proposition}[Separation from numerical evaluation]\label{prop:evaluation}
For the critical mean-chain network
\(F_{1,L}(m^{-1}\sum_jx_j)\), every always-correct algorithm that
returns its activation with absolute error at most \(2^{-20L}\)
must probe all \(m\ge n/2\) relevant inputs. Even an evaluator with
success probability at least \(2/3\) on every context requires
\(m/3\) probes in expectation under uniform inputs. Exact output
sampling has polylogarithmic expected online bit work.
\end{proposition}

